% =====================================================================
% DROP-IN REPLACEMENT for sections 1-2 of changing_networked_mind.tex
% (currently lines 121-214). All citation keys assume the bib entries
% in `bib_entries_to_add.bib` are merged into refs.bib first.
%
% Length: ~4 pages of LNCS body text. Matches current §§1-2 footprint.
% =====================================================================

% ======================================================================
\section{What scalar models miss}
\label{sec:background}
% ======================================================================

When the chemical community released phlogiston it did not revise a single
credence. It rewrote a network of commitments in which combustion, calcination,
respiration, and the theory of acids were all load-bearing on a hub that had
to go. Whether one reads this as a Kuhnian paradigm shift \citep{kuhn1962} or
as competition between two rival research programs \citep{holmes2000}, the
structural fact is the same: phlogiston was hypothesised to account for
combustibility, malleability, ductility, colour, odour, flavour, and heat
\citep{boantza2017}, and late-phlogistic theorists (Cavendish, Priestley,
Kirwan, Macquer) proposed cascading revisions that proved mutually
incompatible \citep{blumenthal2017}---because phlogiston was load-bearing for
too much.

The dominant formalism for modelling such communities is network epistemology:
each agent holds a degree of belief, or credence, between $0$ and $1$, in a
fixed proposition, updates by Bayes, and shares evidence or posteriors with
network neighbours \citep{zollman2010,oconnor2019,polygraphs2024}. This
programme has produced striking results. Zollman's effect---that sparsely
connected communities often outperform densely connected ones, because slow
information flow preserves transient diversity---is the formal vindication of
the old intuition that a field benefits when some of its members stubbornly
resist the consensus. Subsequent work has shown that polarization can emerge
even when agents share evidence, if they disagree on how to interpret it
\citep{michelini2023}, and a parallel literature on the division of cognitive
labour shows that a community hedges its bets best not when its scientists
reason as dutiful truth-seekers but when they chase reward
\citep{kitcher1990,strevens2003}. The PolyGraphs reflection
\citep{polygraphs2024} frames the project explicitly normatively---``how ought
we rationally to form opinions''---and reads persistent disagreement as
``ignorance on the part of the community.''

This framing cannot represent the phlogiston case. The puzzle is not that the
community failed to converge on a single credence: it is that the commitments
were entangled, that revising any one of them forced a re-reading of the
others, and that the disagreement was about which entanglements to keep, not
about a number in $[0,1]$. The closest formal precedent is Thagard's ECHO,
which treats the phlogiston/oxygen debate as constraint satisfaction over a
coherence network with excitatory and inhibitory links between hypotheses and
evidence \citep{thagard1989,thagard1990}. ECHO captures the structural
intuition we need; it does not give us the directed causal dependencies, the
carry-over cost of mind-change, or the asymmetry between expanding and
reducing the network that the next sections derive.

Two more recent lines have begun to address parts of what is missing.
Active-inference accounts of epistemic communities supply an endogenous
precision-based mechanism for lock-in: agents ``tend to sample information in
order to justify their own view of reality, which eventually leads to them
having a high degree of certainty about their beliefs,'' and once this
certainty is reached ``it becomes very difficult to get them to change their
views'' \citep{albarracin2022}. The same author and colleagues have extended
the framework to shared expectations and group intentionality
\citep{albarracin2024protentions}, while a parallel line gives a formal route
from individual to group-level generative models that warns against naive
aggregation: the mapping from individual to collective active inference is
``non-trivial,'' with preference conflicts reducing group precision below
single-agent levels \citep{friston2023federated,tschantz2025}. And the
variational cost of mind-change has been formalized at the single-agent level
as a motivated decision that trades the informational cost of revision against
the utility of holding a belief \citep{hyland2025}, with consequences
including confirmation bias and attitude polarization.

What none of these does is treat \emph{belief structure itself} as the object
of inference. They retain a fixed menu of hypotheses or a fixed coupling
graph, and the dynamics they study live in how the topology of communication
interacts with movement \emph{within} that menu. What no one studies is a
community whose menu is itself the thing being learned, revised, and
transmitted---a community whose paradigm is a directed graph of
interdependent commitments, whose dependencies change as the community
expands and prunes them, and whose dynamics on that changing structure
generate the lag, the lock-in, and the eventual reorganization. The room is
nearly empty, and the reason it is empty is that there has been no cheap,
principled mechanism for ``learning what the model even is'' that could be
run across a population.

% ======================================================================
\section{Why paradigms are networks under structure learning}
\label{sec:hidden}
% ======================================================================

Return to Kuhn's deepest observation, which the scalar models gesture at but
cannot embody: a paradigm shapes what a scientist can perceive, what she
chooses to investigate, and which alternatives she will entertain---three
faces, he insists, of a single fact \citep{kuhn1962}. The reason a scalar
credence over a fixed menu cannot carry this is that a menu has no hidden
interior. It can represent which theory is favoured; it cannot represent the
dependencies that make one theory \emph{cost} more to abandon than another,
nor the way a commitment at the centre of a web of inference quietly
determines what counts, at the periphery, as an anomaly worth noticing.

The pivot that forces everything that follows is that the structure of
reality is itself hidden. A scientist confronting a new domain is not given
the right variables and is certainly not given the dependencies among them.
She does not know, in advance, whether two observed regularities share a
hidden common cause, are directly and lawfully coupled, or merely coincide.
The generative model whose structure she must eventually come to hold is, at
the outset, unknown in its very shape. For an active-inference reader this
move is natural: just as agents must infer hidden states they cannot directly
observe, they must also infer the hidden \emph{dependencies among} those
states. The cognitive-science literature on theory acquisition reaches the
same conclusion from a different starting point. Ullman and Tenenbaum
\citep{ullman2020} argue that conceptual development is best understood as
``resource-constrained, hierarchical Bayesian program induction with
different primitives and constraints,'' with theory learning operating as
``stochastic search at two levels of abstraction---an outer loop in the space
of theories, and an inner loop in the space of explanations or models
generated by each theory'' \citep{kemp2008form,ullmangoodman2012}.
``Symbolic representations,'' they emphasise, ``are better suited to
capturing children's intuitive theories but give rise to a harder learning
problem, which can only be solved by exploratory search.''

This is the framing that lets us name the second-debt our contribution pays.
The standard network-epistemology models keep their communities from
collapsing onto the first locally attractive theory by equipping their agents
with an exploration term: an $\epsilon$ of an $\epsilon$-greedy rule, a
temperature, a switching cost, an exogenous prior strength
\citep{zollman2010,oconnor2019}. The diversity that protects the community is
bought with a free parameter. Once we treat the structure of the model as
itself an object of inference, the exploration term is no longer free: it
is the epistemic value of expanding the model to entertain a commitment not
yet in it, which active inference supplies through expected free energy
\citep{friston2015,friston2017curiosity,millidge2021}. \emph{To explore is to
expand.}

The complementary direction is no less principled. Active inference treats
structure learning as the dual of state inference through two complementary
operations: pruning structure the data do not support, and proposing
structure the data demand \citep{friston2017curiosity,smith2020structure}.
The exact half is Bayesian model reduction: because a reduced model shares
its likelihood with the richer model and differs only in its prior, the
change in evidence under any pruning is available in closed form from the
already-inverted full posterior \citep{friston2018bmr}. The mechanism has
been operationalized in active-inference agents and shown to ``operate in the
absence of further sensory experience\dots redundant parameters are pruned,
and agents form more precise representations,'' running at ``the slowest
timescale'' of a documented hierarchy in which structure learning ``proceeds
by redistributing the products of learning to minimise model complexity''
\citep{smith2022plos}. This is the timescale, and the mechanism, on which a
community decides which of its commitments to keep.

What is new here is to make these two operations the actual object of a
multi-agent dynamics. Existing applications of BMR have run on independent
single agents \citep{smith2022plos} or used reduced priors within a single
agent's sensorimotor loop \citep{priorelli2023}, and the canonical
parametric-empirical-Bayes formulation defines ``group'' as between-subject
neuroimaging rather than as interacting agents \citep{friston2016peb}. The
recent collective-intelligence work that does live at the multi-agent scale
proposes shared generative models but does not operationalize BMR
\citep{friston2024ci,kaufmann2021}. Treating a paradigm as the object the
community is jointly reducing and expanding is therefore an open move, and
it is the move this paper makes.

\paragraph{Two fields, one operator.} The single non-obvious step the rest of
the paper turns on is that the same propagation operator
$\Top = (I-A)^{-1}$ generates two distinct fields on the network. Acting on
the unit source it returns the conservatism $\kappa=\Top\bone$, the
carry-over cost of revising a commitment: a node is costly to revise because
much depends on it. Acting on a utility source $u$ it returns the conviction
field $U=\Top u$, the intrinsic value of holding a belief, propagated through
the same dependencies that carry the cost. The two fields are linearly
independent: a commitment can be central yet value-neutral, cheap yet
cherished, structurally protected yet wanted false. This is the formal
counterpart of the empirical posture of the Causal Attitude Network model,
which treats attitudes as networks of evaluative reactions in which
connectivity is attitude strength \citep{dalege2016}---now realised on the
same directed graph that carries the cost of revision. We argue in
Sections~\ref{sec:bmebmr}--\ref{sec:both} that lock-in, motivated
persistence, Kuhnian transitions, and exploration-as-expansion all fall out
of how these two fields interact.

\paragraph{What the rest of the paper does.} The construction is built to
satisfy six things and the rest of the paper checks them off:
that a paradigm be a structured object with a dense core and a sparse belt;
that revising a commitment incur an aggregate carry-over cost over its
descendants; that conservatism be derived from structural position rather
than assigned by a dial; that structure learning be principled---Bayesian
model reduction over a continuous precision space whose corners are the
discrete topologies; that conviction be a second field on the same network,
decoupled from conservatism and written into the model by the same expansion
that writes a new belief; and that the population signature be a graded
reorganization, or a stall, rather than a single thresholded flip. The
active-inference machinery enters unmodified throughout. Section~\ref{sec:results}
reports how each desideratum fares once the model is actually run.
